
\documentclass[11pt, a4paper]{article}

% Language
\usepackage[english]{babel} 				%option for paper in german
% !TeX spellcheck = en_US

% Font
\usepackage[utf8]{inputenc}					%package to set font encoding
\usepackage[T1]{fontenc}
\usepackage{fontawesome}					
%\usepackage[default]{lato}					%font style
\usepackage{mathpazo}						%font style
\usepackage{eurosym} 						%package to get eurosymbol per \euro
\usepackage{csquotes} 						%flexible quotation marks
\usepackage[super]{nth}						%superscripts for numbers 2nd, 3rd etc
% \usepackage{sansmath}\sansmath				%Turning on sans-serif math mode

% Page Formatting
\setlength{\textheight}{27cm}				%text height
\setlength{\textwidth}{18cm}				%text width
\setlength{\oddsidemargin}{-1cm}			%page margin
\setlength{\evensidemargin}{-1cm} 			%page margin
\setlength{\topmargin}{-3cm}				%top margin
\setlength{\footnotesep}{.5cm}				%space between footnote and text
\usepackage[singlespacing]{setspace} 		%package to adjust line spacing
\setlength{\parindent}{0em}					%set paragraph indentation
\usepackage{pdflscape} 						%package to rotate pages
\usepackage{layouts}						%package to show page formatting settings and convert into different units of measurement

% Bibliography
\usepackage[backend=biber,citestyle=authoryear-comp,bibstyle=authoryear,natbib=true,doi=false,isbn=false,url=false,eprint=false,giveninits=true, maxnames=10, maxbibnames=10,uniquename=false,uniquelist=false]{biblatex}
\addbibresource{~/OneDrive - University of Bristol/10 Literature/Archived/Personal.bib}	%path to library
\AtBeginBibliography{
	\renewcommand*{\mkbibnamefamily}[1]{\textsc{#1}} 			%family names in small caps
}
\renewbibmacro{in:}{}											%No "In: " before journal name
\renewbibmacro*{volume+number+eid}{%							%Volume in brackets
	\printfield{volume}%
	%  \setunit*{\adddot}										% DELETED
	\setunit*{\addnbspace}										% NEW (optional); there's also \addnbthinspace
	\printfield{number}%
	\setunit{\addcomma\space}%
	\printfield{eid}}
\DeclareFieldFormat[article]{number}{\mkbibparens{#1}}


% Hyperlinks
\usepackage{hyperref}
\hypersetup{
	colorlinks,
	linkcolor={red!50!black},
	citecolor={blue!50!black},
	urlcolor={blue!80!black}
}

% Floats
\usepackage{float} 							%allowing for H float specifier
\usepackage{caption}						%define caption font
\captionsetup{labelsep=endash, justification=centering,font=bf,labelfont=sc, skip=5pt, position=top, singlelinecheck=false, parskip=5pt}
\usepackage{placeins}						%package to allow for float barriers
\setlength{\intextsep}{1cm}					%Specify distance of floats to text
\setlength{\floatsep}{1cm} 					%Specify distance between two floats

% Graphics
\usepackage{tikz} 	%package for creating graphics
\usetikzlibrary{calc}


% Tables
\usepackage{booktabs} 						%package for nice tables
\usepackage{multirow} 						%package to combine cells in tables
\usepackage{longtable}						%package to allow for long tables
\usepackage{makecell}						%package to allow line breaks in cells
\usepackage{array}							%package to format tabular environments
%\captionsetup[table]{belowskip=10pt}		%distance text caption
%\captionsetup[longtable]{aboveskip=10pt}	%distance text caption



% Figures
%\captionsetup[figure]{aboveskip=10pt}		%distance text caption
\usepackage{graphicx} 						%package to include graphics
\usepackage{svg}							%include svg images
\usepackage{subfigure}						%package for subfloats

% Math
\usepackage{amsfonts,amssymb, amsthm, mathtools} 	%packages to facilitate writing math
\newtheorem{theorem}{Proposition}					%define "Proposition"-environment
\renewcommand\qedsymbol{$\blacksquare$}				%define qed-symbol
\DeclareMathOperator{\card}{card}
\DeclareMathOperator{\eur}{\text{\euro}}

% Other
\usepackage{comment} 	%package to allow for commenting out entire sections
\usepackage{soul} 		%package to highlight sections
\usepackage{color}		%colorcoding
\newcommand{\hlcyan}[1]{{\sethlcolor{cyan}\hl{#1}}}
\usepackage{enumitem}	%package to allow for enumeration/itemization
\usepackage{titling}	%control package for \maketitle

% JEL/Keywords
\providecommand{\keywords}[1]{\textbf{Keywords:} #1}
\providecommand{\JEL}[1]{\textbf{JEL-Codes:} #1}

%Format Section Headers
\usepackage[explicit]{titlesec} %package for title formatting
\titleformat{\section}{\normalfont\bfseries\large\raggedright\uppercase}{\thesection}{1em}{\large #1}
\titleformat{\subsection}{\normalfont\it\large}{\thesubsection}{1em}{\large #1}
\renewcommand{\thesection}{\arabic{section}}
\renewcommand{\thesubsection}{\thesection.\arabic{subsection}}

%Format Paragraph Headers
%Format Paragraph Headers
\makeatletter
\renewcommand\subparagraph{%
	\@startsection {subparagraph}{5}{\z@ }{0ex \@plus 1ex
		\@minus .2ex}{-1em}{\normalfont \normalsize \it }}%
\makeatother

%Footnotes/Endnotes
\usepackage{endnotes}
%\let\footnote=\endnote

\newcommand{\tabitem}{~~\llap{\textbullet}~~}


\begin{document}
\onehalfspacing
\setlength{\parskip}{.5cm}

%----------------------------------------------------------------------------------------
%	HEADER SECTION
%----------------------------------------------------------------------------------------
\begin{center}
{\LARGE \centering \textbf{Paul Hufe}}\\
{\small \faMousePointer~\href{https://www.paulhufe.net}{www.paulhufe.net}~~~~~~\faTwitter~\href{https://twitter.com/paulhufe?lang=de}{paulhufe}}
\end{center}

\vspace{.5cm}

\section{Courses}

\begin{tabular}{@{}p{4.5cm}@{}p{4.5cm}@{}p{4.5cm}@{}p{4.5cm}}
	2023&\multicolumn{2}{@{}p{9cm}}{Applied Economics (M.Sc.)}& Unit Co-Director\\
	2022--2023&\multicolumn{2}{@{}p{9cm}}{Applied Economic Dissertations (B.Sc.)}& Co-Lecturer\\
	2022&\multicolumn{2}{@{}p{9cm}}{The Economy (B.Sc.)}& Unit Co-Director\\
	2021&\multicolumn{2}{@{}p{9cm}}{Current Economic Problems (B.Sc.)}& Co-Lecturer\\
	2018--2020&\multicolumn{2}{@{}p{9cm}}{Advanced Topics on Inequality (M.Sc.)}& Teaching Assistant\\
	2015--2017&\multicolumn{2}{@{}p{9cm}}{Equality of Opportunity (B.Sc.)}& Unit Director\\
	2012&\multicolumn{2}{@{}p{9cm}}{Advanced International Macroeconomics (M.Sc.)}& Teaching Assistant\\
	---&\multicolumn{2}{@{}p{9cm}}{6 Bachelor and 17 Master theses}& Supervisor\\
\end{tabular}


\section{Teaching Interests} 
\begin{enumerate}
\item Applied Economics (all levels),
\item Public Economics (all levels),
\item Economic Principles (first-year undergraduate),
\item Topics courses on Fairness and Inequality, Gender Economics, Education Economics (third-year undergraduate and above),
\item General interest in interdisciplinary teaching, e.g., the intersection of Economics with Philosophy, Sociology, or Psychology.
\end{enumerate}

\section{Teaching Philosophy}
My teaching philosophy rests on three main pillars. 

First, economics suffers from a significant lack of diversity, particularly in terms of the under-representation of women, ethnic minorities, and individuals from disadvantaged socio-economic backgrounds. This lack of diversity is likely to have important repercussions since we risk missing out on valuable contributions and insights that could help address societal challenges. Therefore, I am firmly committed to the adoption of teaching practices that support under-represented groups to view themselves as part of the economics community. For example, to the extent possible, I try to base my teaching on material from a diverse set of authors and choose teaching examples that represent the wider population. Furthermore, I always provide communication channels that allow to preserve anonymity and thereby reduce participation barriers for all students regardless of their background.

Second, I view economics as a social science and also believe that it has to be taught as such. Therefore, I tailor my teaching material to speak to the societal issues that students are passionate about. For example, even in my first-year class on principles of economics, I try to incorporate applied work from the current research frontier to illustrate how fundamental economic concepts can be used to analyze and understand the societal challenges of our times.


Third, I strive to adopt teaching and assessment practices that decompartmentalize the course material and foster comprehensive student learning. For example, I open most of my classes with multiple choice quizzes that do not only cover the material from the previous week but cover all course content to date. This mode of testing nudges students to stay abreast with all course material and to engage in spaced repetitions and interleaved learning practices. Furthermore, I often rely on essay-based assessments that span a wide range of topics and allow me to test students' ability to connect various elements of the course material. 

\end{document}

